\Needspace{12\baselineskip}
\subsection{Nullstellen und Linearfaktoren}
\label{subsec:real_functions_polynomial_zeros}
Um Nullstellen $x_i$ eines Polynoms zu finden, setzen wir seinen Funktionswert
gleich null:
\[
  p(x_i)=\sum_{k=0}^{n}a_kx_i^k\overset{!}{=}0.
\]
Wie lösen wir diese Gleichung? Für $n \leq 2$ haben wir die Lösungsmethoden
schon kennengelernt. Mit steigendem Grad wird der direkte Weg aber zunehmend
schwieriger:
\par\medskip
{\centering
  \renewcommand{\arraystretch}{1.2}
  \begin{tabular}{@{}>{\centering\arraybackslash}m{1.45cm}|m{4.4cm}|m{8.2cm}|>{\centering\arraybackslash}m{0.7cm}@{}}
    \textbf{Grad} & \textbf{Gleichungsart}
                  & \textbf{Methode zur Nullstellenbestimmung} & \\ \hline
    $n=1$ & lineare Gleichung & Umstellen
          & \raisebox{-.25\height}{\CachedTikzFigure{content/resources/shared/tikz/smiley_happy.tex}}\\ \hline
    $n=2$ & quadratische Gleichung & ABC-Formel
          & \raisebox{-.25\height}{\CachedTikzFigure{content/resources/shared/tikz/smiley_happy.tex}}\\ \hline
    $n=3$ & kubische Gleichung & Formel von Cardano (aufwendig)
          & \raisebox{-.25\height}{\CachedTikzFigure{content/resources/shared/tikz/smiley_sad.tex}}\\ \hline
    $n=4$ & & Formel von Ferrari (aufwendig)
          & \raisebox{-.25\height}{\CachedTikzFigure{content/resources/shared/tikz/smiley_sad.tex}}\\ \hline
    $n\ge5$ & & keine allgemeine Wurzelformel
            & \raisebox{-.25\height}{\CachedTikzFigure{content/resources/chapters/real_functions/tikz/functions_root_method_unavailable.tex}}\\ \hline
  \end{tabular}\par}

\medskip
Oberhalb von $n=2$ gibt es entweder keine oder nur sehr komplizierte allgemeine Lösungsformeln.
Daher wählen wir in diesem Fällen eine andere Strategie:

\begin{TheoremBox}{Nullstelle und Linearfaktor}{real_functions_factor_theorem}
  Hat ein Polynom $p$ vom Grad $n$ die Nullstelle $x_i$, lässt sich der Linearfaktor
  $(x-x_i)$ abspalten.
  \[
    p(x_i)=0
    \quad\Longleftrightarrow\quad
    p(x) = (x-x_i) \cdot q(x)
  \]
  Für das Polynom $q(x)$ gilt:
  \begin{itemize}[nosep,leftmargin=*]
    \item $q$ hat den Grad $n-1$.
    \item Jede Nullstelle von $p(x)$ außer der abgespaltenen Nullstelle $x_i$ ist auch Nullstelle von $q$
  \end{itemize}
\end{TheoremBox}

Um den verbleibenden Faktor zu finden, stellen wir die Gleichung aus
Satz~\ref{thm:real_functions_factor_theorem} nach $q(x)$ um:
\[
  q(x)=\frac{p(x)}{x-x_i}.
\]
Wir bestimmen dieses Polynom $q(x)$ in dem wir den Bruchs mittels \emph{Polynomdivision}
\EnglishTerm{polynomial division} auswerten.

\begin{ExampleBox}{Eine bekannte Nullstelle nutzen}{real_functions_polynomial_division}
  Für $p(x)=x^3-x^2-4x+4$ erkennt man $p(1)=0$.
  Die Polynomdivision durch $x-1$ verringert den Grad.
  \begin{ClozeArea}[34]{11}
    Nacheinander wird jeweils der höchste Term durch $x$ geteilt und das
    zugehörige Vielfache von $x-1$ abgezogen:
    \[
      \begin{array}{r@{\;}r@{\;}r@{\;}r@{\qquad}l}
        x^3&-x^2&-4x&+4&:(x-1)=x^2-4\\
        -(x^3&-x^2)&&&\\ \cline{1-2}
        &&-4x&+4&\\
        &&-(-4x&+4)&\\ \cline{3-4}
        &&&0&\text{Rest}
      \end{array}
    \]
  \end{ClozeArea}
  Der verbleibende Faktor $q(x)=\Cloze{3cm}{x^2-4}$ ist quadratisch.
  Seine Nullstellen bestimmen wir mit der ABC-Formel:
  \begin{ClozeArea}[34]{4}
    \[
      \begin{aligned}
        q(x)=x^2-4&\overset{!}{=}0,\qquad a=1,\ b=0,\ c=-4,\\
        x_{1,2}&=\frac{-b\pm\sqrt{b^2-4ac}}{2a}
        =\frac{0\pm\sqrt{0^2-4\cdot1\cdot(-4)}}{2\cdot1}
        =\pm2.
      \end{aligned}
    \]
  \end{ClozeArea}
  $q$ hat damit alle Nullstellen von $p$ außer der abgespaltenen bei $x=1$.
  Zusammen mit dieser ergeben sich die Nullstellen und die Zerlegung von $p$:
  \[
    \mathcal{N}_p=\{-2,1,2\},\qquad
    p(x)=(x-1)(x-2)(x+2).
  \]
\end{ExampleBox}

\par\Needspace{28\baselineskip}
Dieses Prinzip verallgemeinern wir im folgenden Ablaufdiagramm: Bei Polynomen
höheren Grades spalten wir Linearfaktoren ab, bis ein quadratischer Faktor übrig bleibt.
\par\smallskip
{\centering
  \CachedTikzFigure{content/resources/chapters/real_functions/tikz/functions_polynomial_zero_workflow.tex}\par}
\smallskip

\par\noindent\begin{minipage}{\linewidth}
\begin{DefinitionBox}{Vollständige Linearfaktorzerlegung}{real_functions_multiplicity}
  Eine Darstellung eines Polynoms $p$ vom Grad $n\ge1$ in der Form
  \[
    p(x)=a_n\prod_{i=1}^{r}(x-x_i)^{k_i},
    \qquad \sum_{i=1}^{r}k_i=n,
  \]
  heißt \emph{vollständige Linearfaktorzerlegung}
  \EnglishTerm{complete factorization into linear factors}. Dabei gilt:
  \begin{itemize}[nosep,leftmargin=*]
    \item $a_n\ne0$ ist der Koeffizient von $x^n$.
    \item $x_1,\ldots,x_r\in\R$ sind paarweise verschiedene Nullstellen von $p$.
    \item Die positiven ganzzahligen Exponenten $k_i$ heißen
      \emph{Vielfachheiten} \EnglishTerm{multiplicities}:
      Der Faktor $(x-x_i)$ tritt genau $k_i$-mal auf;
      $x_i$ ist eine \emph{$k_i$-fache Nullstelle}.
  \end{itemize}
\end{DefinitionBox}
In Beispiel~\ref{ex:real_functions_polynomial_division} sind die drei Nullstellen
$-2$, $1$ und $2$ jeweils einfach: $k_1=k_2=k_3=1$.
\par\smallskip
\begin{ExampleBox}[before skip=0.3\baselineskip]{Schneiden oder Berühren?}{real_functions_multiplicity_graph}
  Bei $p(x)=(x+1)(x-1)^2$ ist $x=-1$ eine einfache Nullstelle ($k=1$),
  $x=1$ eine doppelte Nullstelle ($k=2$).
  \par\smallskip
  \begin{minipage}[c]{0.46\linewidth}
    Bei $x=-1$ wechselt $x+1$ sein Vorzeichen: Der Graph
    $\Cloze{3cm}{\text{schneidet}}$ die $x$-Achse.

    Bei $x=1$ bleibt $(x-1)^2$ auf beiden Seiten positiv: Der Graph
    $\Cloze{3cm}{\text{berührt}}$ die $x$-Achse.
  \end{minipage}\hfill
  \begin{minipage}[c]{0.50\linewidth}
    \centering
    \CachedTikzFigure{content/resources/chapters/real_functions/tikz/functions_multiplicity.tex}
  \end{minipage}
\end{ExampleBox}
\end{minipage}

\begin{KeyPointBox}{Vielfachheit und Vorzeichenwechsel}{real_functions_multiplicity_sign}
  An einer reellen Nullstelle eines Polynoms gilt:
  \[
    \begin{array}{c@{\qquad}l}
      k\text{ ungerade}&\text{Vorzeichenwechsel: Achse wird geschnitten},\\
      k\text{ gerade}&\text{kein Vorzeichenwechsel: Achse wird berührt}.
    \end{array}
  \]
  Eine vollständige Linearfaktorzerlegung über $\R$ ist nicht immer möglich:
  $x^2+1$ hat beispielsweise keine reelle Nullstelle.
\end{KeyPointBox}
